\documentclass[11pt]{article}

\usepackage[utf8]{inputenc}
\usepackage[T1]{fontenc}
\usepackage{times}
\usepackage[margin=1in]{geometry}
\usepackage{amsmath,amssymb}
\usepackage{booktabs}
\usepackage{multirow}
\usepackage{graphicx}
\usepackage{hyperref}
\usepackage{url}
\usepackage{xcolor}
\usepackage{caption}
\usepackage{subcaption}
\usepackage{orcidlink}

\hypersetup{
  colorlinks=true,
  linkcolor=blue!50!black,
  citecolor=blue!50!black,
  urlcolor=blue!50!black,
}

\title{\bf SGSzero: Training a Compact Sanguosha Agent\\ via Privileged Rule-Based Supervision}

\author{
Evan Wen\,\orcidlink{0009-0005-7254-8002}\\
Independent Researcher
}

\date{}

\begin{document}

\maketitle

\begin{abstract}
A game agent that runs on a player's own machine has to decide from the information that player can see, and it has to decide fast enough to keep up with the game. Both constraints are easy to state and awkward to satisfy together, because the information that makes training easy is usually the information that deployment removes.

We ran into this directly. Our handcrafted rule engine scores every card in hand by how useful it is. Handing those scores to a small Transformer as input produced a model that reached 83.9\% accuracy on held-out data and 35.6\% the moment the scores were taken away: it had learned to read the number instead of the game. Distilling the scores through an auxiliary head, and discarding that head at deployment, recovers the gap---80.9\% on the same sub-task, from 1.76M parameters, on a CPU.

Three things we did not expect. Hyper-parameter search was useless: seven configurations all landed within 72--73\%, while a 7.6\% increase in data moved accuracy by 5.3 points. Almost all remaining errors are low-confidence boundary calls rather than systematic misunderstanding---91\% of them. And the aggregate metric hid a subgroup that had collapsed to near-baseline, hands with no attack card, which the auxiliary head repairs by 19.6 points.
\end{abstract}

\section{Introduction}

Card game AI has produced landmark results: AlphaGo and AlphaZero for perfect-information board games \cite{silver2016,silver2018}, Libratus and Pluribus for poker \cite{brown2018,brown2019}, Suphx for mahjong \cite{li2020suphx}, and DouZero for DouDizhu \cite{zha2021douzero}. Each of these required substantial compute and was evaluated in a setting where the agent observes the full game state internally, regardless of what a human player could see.

Two practical constraints are usually left out of this line of work.

\textbf{Constraint 1: human-observable information only.} In a deployed assistant or a training tool for human players, the agent cannot access the opponent's hand, the deck order, or the engine's internal evaluation of card utility. It must decide from exactly the information a human player has.

\textbf{Constraint 2: consumer-hardware latency.} If the agent is to run alongside a live game client---or on a phone---it cannot call a large language model. Inference must complete in milliseconds on a CPU.

These constraints interact in a way that creates a genuine learning problem. A rule-based engine, written by a domain expert, \emph{does} have access to internal valuations: it scores each card in hand by how useful it would be. This information is highly informative for the decision, but it is unavailable at deployment. The naive approach---feed the valuation into the model as an input feature---produces a model that learns to read the number instead of learning the game. It performs well in-distribution and collapses when the feature is absent.

We encountered exactly this failure. A model trained with valuations as input reached 83.9\% accuracy on held-out data but dropped to 35.6\% when the valuation feature was removed, as it must be in deployment.

Our contribution is a practical recipe that resolves this, together with a systematic empirical study of what actually drives performance in this regime. Specifically:

The contributions are as follows. We formulate the Sanguosha decision problem under a human-observable information constraint and release a state encoding that matches what a player actually sees (Section~\ref{sec:problem}). We then separate two kinds of privileged signal that are often conflated: closing an \emph{observability} gap, as perfect-information distillation does \cite{yang2022perfectdou}, and closing a \emph{knowledge} gap, as here. The distinction is not cosmetic. Revealing hidden state cannot close a knowledge gap, because a network that can see every card still has to learn what a card is worth. Our method applies LUPI \cite{vapnik2009,vapnik2015} as an auxiliary card-value head trained on the engine's valuations and discarded before deployment, which recovers most of the gap (78.0--78.4\% to 80.8--81.8\%) at no inference cost (Sections~\ref{sec:method}--\ref{sec:experiments}). Finally, we report the negative and diagnostic results in some detail, because in this regime they were more informative than the headline number (Section~\ref{sec:analysis}).

\section{Related Work}

\paragraph{Card and board game AI.} AlphaZero \cite{silver2018} combines Monte Carlo tree search with a policy-value network trained purely by self-play. DouZero \cite{zha2021douzero} trains a Deep Monte-Carlo value network for DouDizhu with a compact representation and no human data. Suphx \cite{li2020suphx} reaches superhuman mahjong through self-play with global reward prediction. Our setting differs in that the game involves chained skill effects rather than a fixed action vocabulary, and that we explicitly forbid information unavailable to a human player.

\paragraph{Training with information unavailable at deployment.} The closest line of work to ours is \emph{PerfectDou} \cite{yang2022perfectdou}, which introduces \emph{perfect information distillation}: agents are trained as if the game were fully observable---the network sees every player's hand---and the resulting policy is then executed under partial observability. This is a \emph{perfect-training, imperfect-execution} framework, and it shares our high-level philosophy that privileged signals available only during training can substantially improve a deployable policy.

The two settings differ in what the privileged signal \emph{is}, and the difference matters. In PerfectDou the privileged signal is a part of the game state that objectively exists but is hidden from the player: the opponent's cards. Distillation here removes an \emph{observability} gap. In our setting the privileged signal is the rule engine's scalar valuation of each card---a quantity that does not exist anywhere in the game state and is produced by an external symbolic program. What we remove is a \emph{knowledge} gap, not an observability gap. A perfect-information formulation cannot express it: even with all cards revealed, a network still has to learn what a card is worth.

Consequently the two approaches are complementary rather than competing. PerfectDou could in principle be layered on top of our method (observe everything \emph{and} receive valuations), and our contribution is orthogonal: how to inject a symbolic expert's internal valuation into a compact neural policy without letting that valuation dominate the learned representation.

\paragraph{Privileged information.} Learning Using Privileged Information (LUPI) was introduced by Vapnik and Vashist \cite{vapnik2009} and extended in \cite{vapnik2015}: a teacher has access to additional information $x^{*}$ at training time that is unavailable at test time. Distillation-style variants transfer such knowledge through softened teacher outputs. A known pitfall, discussed in the LUPI literature, is that if the privileged information dominates the shared representation, the student under-uses the primary features and generalizes poorly once the privileged channel is removed. Our results reproduce this failure mode and show that a low-weight auxiliary head avoids it. % TODO(author): insert the two specific LUPI/distillation references here before submission.

\paragraph{Compact models for structured decisions.} Transformers \cite{vaswani2017} over tokenized state sequences have become a standard way to encode heterogeneous game state. Our encoder is deliberately small (1.76M parameters) to satisfy the deployment constraint, following the broader argument that domain-specialized compact models can outperform far larger general models on narrow tasks \cite{gunasekar2023}.

\section{Problem Setup}
\label{sec:problem}

\subsection{The Game}

\emph{Sanguosha} (Legends of the Three Kingdoms) is a multi-player Chinese card game. Each player controls a hero with a unique skill, holds a hand of cards, and acts in turn. Effects chain: playing a card may trigger skills, which may trigger further skills, and resolution order matters. We study the two-player \emph{DouDizhu-style} variant implemented by the open-source client we use as an environment, in which identities are public and only the top of the deck is hidden.

We focus on \textbf{Task A}: at the play phase, decide whether to play a card or pass. This is the highest-frequency decision in the game and the prerequisite for the harder \emph{which card} task.

\subsection{Human-Observable State}

The input is a token sequence. Table~\ref{tab:state} lists what is included and, critically, what is excluded.

\begin{table}[t]
\centering
\small
\begin{tabular}{@{}p{0.46\linewidth}p{0.46\linewidth}@{}}
\toprule
\textbf{Included (human-observable)} & \textbf{Excluded (engine-internal)} \\
\midrule
Own hand: card name, rank, suit & Card utility score from the rule engine \\
All heroes' skills (public) & Internal event stream \\
All players' HP and max HP & Turn index, seat index \\
All identities (public in this mode) & Cards already played this round \\
Opponents: hero, HP, hand \emph{count}, identity, alive flag & Opponent hand \emph{contents} \\
Deck remaining \emph{count} & Top card of deck \\
Equipments, judgement area, chained/flipped status, marks & \\
Current phase & \\
\bottomrule
\end{tabular}
\caption{State features. The left column is what a human player can observe; the right column is available to the engine but withheld from the model, because a human cannot see it.}
\label{tab:state}
\end{table}

The resulting sequence has a vocabulary of 561 tokens and a 95th-percentile length of 51 tokens. Special structural tokens \texttt{[BOS]}, \texttt{[SELF]}, \texttt{[HAND]}, \texttt{[OTHER]}, \texttt{[EOS]} delimit segments.

\subsection{Privileged Signal}

The rule engine we use as a teacher is a 10,966-line handcrafted policy. For each card in hand it exposes a scalar utility $v \in \{0,\dots,12\}$. This is precisely the kind of information a strong human player forms implicitly and a rule engine computes explicitly---and precisely what is unavailable at deployment.

\section{Method}
\label{sec:method}

\subsection{Encoder}

A token embedding of dimension $d=256$ feeds three Transformer encoder layers (4 heads, feed-forward 512, pre-norm, dropout 0.1). The decision representation is taken from the \texttt{[SELF]} position:

\begin{equation}
h = \mathrm{LayerNorm}\big(\mathrm{Encoder}(x)_{[\texttt{SELF}]}\big), \qquad
\hat{y} = \mathrm{softmax}(W h + b).
\end{equation}

We use the \texttt{[SELF]} position rather than mean pooling deliberately. Mean pooling over a variable-length hand produces representations whose scale depends on hand size, conflating ``how many cards'' with ``what cards''. The \texttt{[SELF]} token attends over the whole sequence and therefore carries an aggregate summary whose scale is independent of hand size. Total parameters: 1.76M.

\subsection{Privileged Value Head}

During training only, we attach an auxiliary head that predicts the engine's utility bucket for each card in hand:

\begin{equation}
\mathcal{L} = \mathcal{L}_{\mathrm{dec}} + \lambda \, \mathcal{L}_{\mathrm{val}},
\qquad
\mathcal{L}_{\mathrm{val}} = \frac{1}{|H|}\sum_{i \in H} \mathrm{CE}\big(\mathrm{MLP}(x_i),\, \mathrm{bin}(v_i)\big),
\end{equation}

where $H$ is the set of card positions in the hand and $v_i$ is the teacher's utility for card $i$, discretized into 13 bins. We use $\lambda = 0.5$.

The head is discarded at deployment. Its purpose is representational: by requiring the encoder to predict card utility from card identity, rank, and suit alone, it forces the encoder to \emph{internalize} a utility function rather than read one off an input feature.

This is the LUPI recipe applied to a rule engine rather than to a human expert. The teacher is symbolic and cheap; the student is neural and fast.

\section{Experiments}
\label{sec:experiments}

\subsection{Data}

Training data is generated by letting the rule engine play against itself in the live game client, logging (state, action) pairs. The latest dataset contains 13,536 records with the full observable feature set, split into 9,470 training and 340 validation records by game ID (no game appears in both splits).

\subsection{Baselines}

We report the majority-class baseline on each validation split, which is the accuracy of always predicting the more frequent label. This is the honest floor: a model that only learns the label prior scores exactly this.

\subsection{Main Results}

\begin{table}[t]
\centering
\small
\begin{tabular}{@{}lcccc@{}}
\toprule
\textbf{Configuration} & \textbf{Input} & \textbf{Val size} & \textbf{Accuracy} & \textbf{macro-F1} \\
\midrule
V5 (naive privileged input) & full state + $v,u$ & 819 & 83.9\% & --- \\
\quad same, at deployment & observable only & --- & \textbf{35.6\%} & --- \\
\midrule
V6 (observable only) & observable & 819 & 78.0--78.4\% & --- \\
V7 (+ privileged head) & observable & 819 & 80.8--81.8\% & --- \\
V8 (+ full fields, + head) & observable & 340 & \textbf{80.6--80.9\%} & 0.756--0.762 \\
\midrule
Majority baseline & --- & 340 & 69.2\% & --- \\
\bottomrule
\end{tabular}
\caption{Accuracy across configurations. V5 shows the failure mode: high in-distribution accuracy collapses when the privileged feature is removed. V7/V8 recover most of the gap while remaining deployable.}
\label{tab:main}
\end{table}

Table~\ref{tab:main} summarizes the progression. The key comparison is between V5 and V7. V5 is \emph{more accurate} in-distribution (83.9\%) but is unusable: it depends on a feature that does not exist at deployment, and removing it costs 48.3 points. V7 gives up 2.1 points of in-distribution accuracy in exchange for being deployable at all.

V8 adds previously missing observable fields (chained/flipped status, marks) and lands at 80.9\%, statistically indistinguishable from V7 despite a smaller validation set, confirming that the new fields entered the representation without displacing anything.

\subsection{Deployment Cost}

The model has 1,760,527 parameters (6.72 MB in fp32). Training 40 epochs takes 70 seconds on a single RTX 4060 Laptop GPU. Inference is a single forward pass over a sequence of at most $\sim$100 tokens and runs on commodity CPUs in milliseconds, comfortably inside a human player's decision window.

\section{Analysis}
\label{sec:analysis}

\subsection{Capacity Is Not the Bottleneck; Data Is}

We swept seven hyper-parameter configurations (width $\in \{128,192,256\}$, depth $\in \{2,3,4\}$, learning rate $\in \{2\text{e-}4, 3\text{e-}4, 1\text{e-}3\}$). Every configuration landed in 72.0--73.0\%. Increasing data from 17,830 to 19,188 records (+7.6\%) under a \emph{fixed} configuration moved accuracy from 73.0\% to 78.3\% (+5.3 points).

\begin{table}[t]
\centering
\small
\begin{tabular}{@{}lcc@{}}
\toprule
\textbf{Configuration} & \textbf{Accuracy} & \textbf{macro-F1} \\
\midrule
$d{=}128$, $L{=}2$, lr $3\text{e-}4$ & 72.3\% & 0.702 \\
$d{=}128$, $L{=}2$, lr $1\text{e-}3$ & 72.0\% & 0.692 \\
$d{=}256$, $L{=}3$, lr $3\text{e-}4$ & 73.0\% & 0.692 \\
$d{=}128$, $L{=}4$, lr $3\text{e-}4$ & 72.0\% & 0.695 \\
$d{=}192$, $L{=}3$, lr $2\text{e-}4$ & 72.8\% & 0.697 \\
\midrule
$d{=}256$, $L{=}3$, \emph{more data} & \textbf{78.3\%} & 0.756 \\
\bottomrule
\end{tabular}
\caption{Hyper-parameter sweep versus data scaling. Model size, depth, and learning rate are irrelevant in this regime; the data change is not.}
\label{tab:hyper}
\end{table}

The practical implication is blunt: in this regime, spending effort on architecture is waste. We tuned seven configurations over several days and gained half a point; a few hours of data collection gained five. The binding constraint is the quantity and quality of logged decisions.

\subsection{Data Quality Dominates Quantity}

A separate comparison shows 11,084 records with the complete feature set matching 29,834 mixed records (77.4\% versus 78.4\%). One clean record is worth roughly 2.7 noisy ones. The reason is distributional: 64\% of the mixed corpus lacked opponent-action features, so the model was trained on a distribution in which the same input sometimes carried information and sometimes did not.

\subsection{Errors Are Boundary Cases, Not Misconceptions}

Of 98 errors, 39 were low-confidence ($p<0.7$), 50 medium ($0.7\le p<0.9$), and only 9 high-confidence ($p\ge 0.9$). Calibration is close to diagonal: predictions made at $\ge 95\%$ confidence are correct 98.0\% of the time, which we did not expect from a model this small and this sparsely trained. In other words, the model is not confused about the game; it is uncertain at genuinely close decisions.

\subsection{A Local Collapse Hidden by the Aggregate}

The headline accuracy conceals a severe weakness. Splitting the validation set by whether the hand contains an attack card:

\begin{table}[t]
\centering
\small
\begin{tabular}{@{}lcccc@{}}
\toprule
\textbf{Subset} & \textbf{$n$} & \textbf{Baseline} & \textbf{No head} & \textbf{With head} \\
\midrule
Holds attack card & 350 & 74.9\% & 82.9\% (+8.0) & 78.7\% (+7.2) \\
No attack card & 103 & 60.2\% & 63.1\% (+2.9) & \textbf{77.3\% (+19.6)} \\
\bottomrule
\end{tabular}
\caption{Subgroup accuracy. Without the privileged head, the model is barely above the majority baseline when it holds no attack card---exactly the situations that require judging card utility. The auxiliary head repairs this.}
\label{tab:subgroup}
\end{table}

When no attack card is held, the decision reduces to whether to spend a trick card, an equipment, or a peach---a pure question of card value. Without privileged supervision the model has no representation of value and falls back on the prior. The auxiliary head lifts this subgroup from +2.9 to +19.6 points over baseline.

An oversampling control confirms the diagnosis: tripling the weight of this subgroup \emph{lowers} accuracy (78.4\% $\to$ 77.9\%). The deficiency is informational, not statistical.

\subsection{Headroom}

Identical token sequences with conflicting labels occur in only 1.9\% of the data, implying a Bayes-optimal ceiling near 99.1\%. With 18,691 distinct states across 18,881 samples, almost every position is unique, so the model must generalize rather than memorize. The gap between 81\% and 99\% is real headroom, not label noise.

\subsection{Validation Protocol}

We ran 13 checks before trusting any number, including: single-batch overfitting (reaches 100\%), shuffled-label control (67.7\% versus 68.4\% baseline, i.e.\ no leakage), padding-mask consistency (difference $6\times10^{-8}$ between single and batched inference), UNK rate (0.0076\%), and game-level train/validation disjointness (zero overlap). We consider the shuffled-label and mask-consistency checks non-negotiable for this kind of pipeline; both are easy to skip and both catch silent bugs.

\section{Limitations}

Our evaluation is restricted to Task A (play or pass) in the two-player DouDizhu-style mode. The \emph{which card} task, response decisions, and the multi-player identity mode are not yet evaluated. Validation sets are small (340--819 records), so subgroup numbers carry wide confidence intervals; the ``no attack card'' result in particular rests on 103 records. The teacher is a handcrafted engine, so the student inherits its strategic ceiling until self-play improvement is added; we regard self-play fine-tuning as the natural next step and the point at which the ``zero human data'' claim becomes fully accurate. Finally, the game client we instrument is a third-party open-source implementation; results may not transfer verbatim to the commercial client.

\section{Conclusion}

The gap between what a rule engine knows and what a player can see turned out to be the hard part of this problem, not the game itself. A rule engine's internal valuations cannot be used at deployment, but they can be used at training time as privileged information---provided they supervise an auxiliary head rather than enter as features. Feeding them in directly produces a model that reads numbers instead of playing the game, and collapses by 48 points at deployment.

Along the way we found that model capacity is irrelevant in this regime, that data cleanliness matters more than volume, that aggregate accuracy conceals a subgroup collapse precisely where card valuation is needed, and that the privileged head repairs it. The resulting model has 1.76M parameters, trains in 70 seconds on a consumer GPU, and runs inference on a CPU.

We release the state encoding, training pipeline, and validation protocol, and we intend to release the self-play improved agent and a benchmark for Sanguosha decision-making.

\section*{Acknowledgment}

We thank the maintainers of the open-source Sanguosha client used as our environment, and the authors of the handcrafted rule engine that served as our privileged teacher.

\begin{thebibliography}{99}

\bibitem{silver2016}
D.~Silver, A.~Huang, C.~J. Maddison, et al.
\newblock Mastering the game of Go with deep neural networks and tree search.
\newblock \emph{Nature}, 529(7587):484--489, 2016.

\bibitem{silver2018}
D.~Silver, T.~Hubert, K.~Schrittwieser, et al.
\newblock A general reinforcement learning algorithm that masters chess, shogi, and Go through self-play.
\newblock \emph{Science}, 362(6419):1140--1144, 2018.

\bibitem{brown2018}
N.~Brown and T.~Sandholm.
\newblock Superhuman AI for heads-up no-limit poker: Libratus beats top professionals.
\newblock \emph{Science}, 359(6374):418--424, 2018.

\bibitem{brown2019}
N.~Brown and T.~Sandholm.
\newblock Superhuman AI for multiplayer poker.
\newblock \emph{Science}, 365(6456):885--890, 2019.

\bibitem{li2020suphx}
J.~Li, S.~Koyamada, Q.~Ye, et al.
\newblock Suphx: Mastering Mahjong with deep reinforcement learning.
\newblock \emph{arXiv:2003.13590}, 2020.

\bibitem{zha2021douzero}
D.~Zha, J.~Xie, W.~Ma, et al.
\newblock DouZero: Mastering DouDizhu with self-play deep reinforcement learning.
\newblock In \emph{Proceedings of the 38th International Conference on Machine Learning (ICML)}, 2021.

\bibitem{yang2022perfectdou}
G.~Yang, M.~Liu, W.~Hong, W.~Zhang, F.~Fang, G.~Zeng, and Y.~Lin.
\newblock PerfectDou: Dominating DouDizhu with perfect information distillation.
\newblock \emph{arXiv:2203.16406}, 2022.

\bibitem{vapnik2009}
V.~Vapnik and A.~Vashist.
\newblock A new learning paradigm: Learning using privileged information.
\newblock \emph{Neural Networks}, 22(5--6):544--557, 2009.

\bibitem{vapnik2015}
V.~Vapnik and R.~Izmailov.
\newblock Learning using privileged information: Similarity control and knowledge transfer.
\newblock \emph{Journal of Machine Learning Research}, 16:2023--2049, 2015.

\bibitem{vaswani2017}
A.~Vaswani, N.~Shazeer, N.~Parmar, et al.
\newblock Attention is all you need.
\newblock In \emph{Advances in Neural Information Processing Systems (NeurIPS)}, 2017.

\bibitem{gunasekar2023}
S.~Gunasekar, Y.~Zhang, J.~Aneja, et al.
\newblock Textbooks are all you need.
\newblock \emph{arXiv:2306.11644}, 2023.

\end{thebibliography}

\end{document}
